\begin{ouverture}[geometrie]{17}{La symétrie, axes droits et obliques}{Partie 4 · Espace et géométrie}
{Axes · Quadrillage · Diagonales · Figures complétées}
\objectif{Reconnaître les axes de symétrie d’une figure.}
\objectif{Construire un symétrique sur quadrillage, axe horizontal ou vertical.}
\objectif{Construire un symétrique le long d’une diagonale du quadrillage.}
\prerequis{Symétrie du CM1 sur axes horizontaux et verticaux.}
\end{ouverture}
\pagestyle{otMain}\essentiel\fontsize{11}{15}\selectfont
\phantomsection\label{cours-17-1}
\essSection{1}{Les acquis, consolidés}
\begin{definition}Une droite est un axe de symétrie d’une figure si le pliage le long de cette droite superpose exactement les deux parties. Le carré a 4 axes, le rectangle 2, le triangle équilatéral 3, le cercle une infinité.\end{definition}
\label{fig-cours-17-1-cours-symetrie}\figureGeo{cours-symetrie}

\begin{savoirfaire}[Compter les carreaux]Axe vertical ou horizontal : le symétrique d’un point se place de l’autre côté, à autant de carreaux de l’axe, sur la même ligne ou la même colonne. Un point de l’axe ne bouge pas.\end{savoirfaire}
\begin{propriete}La symétrie conserve les longueurs, les angles et les aires : le symétrique d’un segment de 5 carreaux mesure 5 carreaux ; celui d’un carré reste un carré superposable.\end{propriete}
\clearpage
\phantomsection\label{cours-17-2}
\essSection{2}{L’axe en diagonale}
\begin{definition}Sur un quadrillage, une diagonale des carreaux peut servir d’axe : c’est la droite qui traverse les carreaux de coin en coin, à 45 degrés.\end{definition}
\label{fig-cours-17-2-cours-diagonale}\figureGeo{cours-diagonale}

\begin{savoirfaire}[Échanger les directions]Avec un axe diagonal, les déplacements s’échangent : un point situé à 2 carreaux à droite de l’axe, en comptant perpendiculairement par les diagonales des carreaux, se reflète à 2 carreaux de l’autre côté. En pratique : un trajet « 3 vers la droite » devient « 3 vers le haut », et réciproquement, quand l’axe monte vers la droite.\end{savoirfaire}
\begin{exemple}Le symétrique du point situé 3 carreaux à droite et 1 carreau au-dessus d’un point de l’axe est le point situé 1 carreau à droite et 3 carreaux au-dessus : les deux nombres s’échangent.\end{exemple}
\begin{attention}Contrôler au calque : plié le long de la diagonale, le quadrillage se superpose à lui-même, et la figure doit tomber sur son image. Un point de l’axe reste immobile, comme toujours.\end{attention}

\clearpage\rubrique{Découverte}
\begin{blocLivre}{ex-17-1}\begin{exercice}[1]{Combien d’axes ?}
Donner le nombre d’axes de symétrie de chaque figure dessinée : carré, rectangle, triangle équilatéral, cercle.
\label{fig-ex-17-1}\figureGeo{sym-1}
\end{exercice}\end{blocLivre}
\begin{blocLivre}{ex-17-2}\begin{exercice}[1]{Axe vertical}
Construire le symétrique de la figure par rapport à l’axe vertical tracé.
\label{fig-ex-17-2}\figureGeo{sym-2}
\end{exercice}\end{blocLivre}
\begin{blocLivre}{ex-17-3}\begin{exercice}[1]{Axe horizontal}
Construire le symétrique du drapeau par rapport à l’axe horizontal.
\label{fig-ex-17-3}\figureGeo{sym-3}
\end{exercice}\end{blocLivre}
\begin{blocLivre}{ex-17-4}\begin{exercice}[1]{Le point sur la diagonale}
L’axe est la diagonale tracée. Placer le symétrique de chacun des trois points marqués.
\label{fig-ex-17-4}\figureGeo{sym-4}
\end{exercice}\end{blocLivre}
\begin{blocLivre}{ex-17-5}\begin{exercice}[1]{Segment et diagonale}
Construire le symétrique du segment [AB] par rapport à la diagonale tracée.
\label{fig-ex-17-5}\figureGeo{sym-5}
\end{exercice}\end{blocLivre}

\rubrique{Entraînement}
\begin{blocLivre}{ex-17-6}\begin{exercice}[2]{La figure oblique}
Compléter la figure pour que la diagonale soit son axe de symétrie.
\label{fig-ex-17-6}\figureGeo{sym-6}
\end{exercice}\end{blocLivre}
\begin{blocLivre}{ex-17-7}\begin{exercice}[2]{Chercher l’erreur}
Noa a tracé le symétrique du triangle par rapport à la diagonale, mais un sommet est faux. Lequel ? Corriger.
\label{fig-ex-17-7}\figureGeo{sym-7}
\end{exercice}\end{blocLivre}
\begin{blocLivre}{ex-17-8}\begin{exercice}[2]{Les lettres au miroir}
Parmi les lettres N, S, Z dessinées, aucune n’a d’axe de symétrie, et pourtant elles semblent « symétriques ». Expliquer.
\label{fig-ex-17-8}\figureGeo{sym-8}
\end{exercice}\end{blocLivre}
\begin{blocLivre}{ex-17-9}\begin{exercice}[2]{Deux axes perpendiculaires}
Compléter la figure pour qu’elle admette l’axe vertical et l’axe horizontal tracés.
\label{fig-ex-17-9}\figureGeo{sym-9}
\end{exercice}\end{blocLivre}
\begin{blocLivre}{ex-17-10}\begin{exercice}[2]{Les longueurs conservées}
Le segment [RS] mesure 4 carreaux le long du quadrillage. Sans mesurer, donner la longueur de son symétrique par rapport à la diagonale, puis contrôler sur la construction.
\label{fig-ex-17-10}\figureGeo{sym-10}
\end{exercice}\end{blocLivre}

\rubrique{Problèmes}
\begin{blocLivre}{ex-17-11}\begin{exercice}[3]{Le logo du club}
Compléter le logo pour que la diagonale tracée soit un axe de symétrie, en coloriant les cases utiles.
\label{fig-ex-17-11}\figureGeo{sym-11}
\end{exercice}\end{blocLivre}
\begin{blocLivre}{ex-17-12}\begin{exercice}[3]{Le pavage malin}
Sur le quadrillage, la frise se poursuit par symétries successives d’axes verticaux régulièrement espacés. Dessiner le motif suivant.
\label{fig-ex-17-12}\figureGeo{sym-12}
\end{exercice}\end{blocLivre}
\begin{blocLivre}{ex-17-13}\begin{exercice}[3]{Axe à trouver}
La figure complète admet un axe de symétrie : le tracer, puis justifier au calque.
\label{fig-ex-17-13}\figureGeo{sym-13}
\end{exercice}\end{blocLivre}
\begin{blocLivre}{ex-17-14}\begin{exercice}[3]{La rosace du vitrail}
Le quart de rosace est dessiné ; les deux axes, vertical et horizontal, sont tracés. Compléter la rosace et dire combien de cases colorées compte la figure finale.
\label{fig-ex-17-14}\figureGeo{sym-14}
\end{exercice}\end{blocLivre}
\begin{blocLivre}{ex-17-15}\begin{exercice}[3]{Vrai ou faux ?}
« Le symétrique d’un carré peut être un rectangle non carré. » « Un point de l’axe est son propre symétrique. » Justifier.
\label{fig-ex-17-15}\figureGeo{sym-15}
\end{exercice}\end{blocLivre}

\clearpage
